\chapter{Noise Robustness Across Molecular Representations and Model Architectures}\label{chap:robustness}
% TODO: Chapter 3 update pending. When it lands, search for `TODO: Ch3 update' in Chapters 1, 4 and 7 and re-check each quoted result. Also re-check the Chapter 3 Additional files and the abstract.
% [notes moved to thesis_notes_moved.md]

% TODO: review, needs manual review
% [CODE INPUT: commented out -- the Chapter 3 paper was not submitted, so it is not referenced]
% \textit{Preamble.} The work in this chapter is adapted from Punt, A., Hanser, T., Werner, S., and Morris, G.M., in preparation for the \textit{Journal of Cheminformatics}. Most of the figures, tables, and text in this chapter are therefore either identical or similar to the content of that manuscript. A signed Statement of Authorship for this jointly authored work is included in the front matter.

\section{Introduction}
% [notes moved to thesis_notes_moved.md]


While there have been many investigations into the effect of noise on molecular property prediction, direction comparisons between how different molecular representations responds to label noise remains limited. The choice of molecular representation certainly affects predictive performance. \citet{Deng2023} evaluated 62,820 models and found that extended-connectivity fingerprints (ECFP) fingerprints frequently outperform learned graph embeddings across models. Whether these conclusions translate to noise robustness is our question. Between different types of label noise, \citet{Heid2023} found no difference in model performance between error distributions with the same mean and standard deviation. We expand on this by considering alternative patterns of noise, namely noise concentrated in structural groups or on a sparse subset of samples.

% TODO: review, needs manual review
This chapter answers the first four questions set out in Chapter~\ref{chap:intro}: how representation and model architecture divide the variance in predictive performance and in noise robustness, what probabilistic models add, whether the type of label noise matters, and whether uncertainty estimates track label noise.

\section{Methods}\label{sec:robustness:methods}
% ---- NOTES: Methods ----
% Takeaway: the study crosses four datasets, six representations, 20 model configurations, seven noise conditions and seven noise levels.
% Paragraph 1 (only if the section gets a lead-in paragraph):
% - Four datasets: QM9 (HOMO--LUMO gap) and three assay datasets (logD, Caco-2, hERG K$_i$).
% - Six representations, 20 model configurations, seven noise conditions, seven noise levels.
% - The 20 configurations are thirteen base models and seven variant models. Name both terms here; the Models subsection defines them.
% - The noise conditions are defined in Section~\ref{sec:noise_injection}. Hyperparameters are in Additional file~\ref{app:noise:sec:hyperparameters_default}.
% ---- END NOTES ----

% CODE INPUT
We used four datasets: QM9, with the HOMO--LUMO gap as the target, and three assay datasets, logD, Caco-2 and hERG K$_i$. On each we fitted six representations and 20 model configurations. The 20 configurations are thirteen base models and seven variant models, both defined under Models below. Every pairing was exposed to seven noise conditions, defined in Section~\ref{sec:noise_injection}, at seven noise levels. Hyperparameters for every model are given in Additional file~\ref{app:noise:sec:hyperparameters_default}.



\subsection{Datasets}

Although we are investigating the effects of label noise stemming from experimental settings, for the purposes of testing we inject controlled amounts of artificial noise into clean labels. Data were extracted from the QM9 molecular property dataset distributed by PyTorch Geometric \citep{Ramakrishnan2014}, which contains 130,831 small organic molecules with pre-computed quantum-mechanical properties obtained from density functional theory (DFT) calculations at the B3LYP/6-31G(2df,p) level of theory. We selected the gap between the highest occupied and the lowest unoccupied molecular orbital (HOMO--LUMO gap) as the prediction target in QM9, as it captures a molecule's electronic excitability, charge transfer capability, and chemical stability \citep{Fediai2023}. We assumed that QM9 labels were free of noise, notwithstanding approximations in the level of theory. Each experiment was replicated ten times with a distinct random seed and a randomly selected subset of $N = 10{,}000$ molecules, split 80/10/10 into 8,000 training, 1,000 validation and 1,000 test molecules. All models were fitted on those 8,000 training molecules, except the Gaussian processes, which were fitted on a random subset of 5,000 of them; the validation and test molecules were unchanged.

We repeated these experiments on three molecular property datasets with experimentally determined endpoints. From the OpenADMET initiative \citep{openadmet}, we used LogD and Caco-2 efflux permeability, with $N = 5{,}039$ and $N = 2{,}161$ respectively. We also selected hERG $K_i$ data from ChEMBL \citep{Zdrazil2023}, following a protocol inspired by \citet{landrum2024}: filter for binding assays, deduplicate by median pChEMBL value, and remove compounds with an inter-assay standard deviation greater than 1.0 log unit, resulting in $N = 1{,}415$ compounds. Each of the three was partitioned into five folds by grouped $k$-fold cross-validation, and each fold's training block gave up a further grouped fifth as a validation set for early stopping. A model therefore fits between 3,194 and 3,364 of the 5,039 LogD molecules, depending on the fold. Unlike QM9's ten replicates, which are ten independent draws from one pool, each assay dataset is a single fixed set in which every molecule is tested exactly once and nothing is redrawn.

All splits and cross-validation folds were grouped by the chirality-blind Bemis--Murcko scaffold, implemented such that each distinct acyclic molecule forms its own group, and groups are filled in random order. Prior to modeling, all molecular structures were sanitized and canonicalized using \texttt{RDKit} to remove invalid valence states and standardize atom and bond typing \citep{rdkit}. Prediction targets were standardized to zero mean and unit variance. Experiments that track uncertainty were run once rather than in replicates.

\subsection{Representations}

In this study, we evaluated six molecular representations: extended-connectivity fingerprints of diameter four (ECFP4), Avalon, Sort \& Slice (SNS), physicochemical descriptors (PDV), ChemBERTa and Molecular Hypergraph Grammar GNNs (MHG-GNNs). The same six were used on QM9 and on the three assay datasets. ECFP4 is a Morgan fingerprint with radius $r = 2$ with $d = 2048$ bits \citep{rogers2010, rdkit}. Sort \& Slice takes the same substructures as ECFP4, but it ranks them by prevalence and selects the top $L$ ($L=1024$) \citep{sns}. We fitted the Sort \& Slice vocabulary on the training molecules alone, refitting it for each QM9 replicate and for each cross-validation fold.  Avalon is a 2,048-bit fingerprint, generated with the Avalon toolkit through RDKit \citep{avalon, rdkit}. The PDV representation is a vector of 200 physicochemical descriptors computed with RDKit's \texttt{MolecularDescriptorCalculator} \citep{rdkit, Cherkasov2014}; the 200 are listed in Additional file~\ref{app:noise:sec:pdv_descriptors}, and the list is fixed in our code rather than taken from RDKit's default descriptor set, which grows between releases. ChemBERTa embeddings are the mean over the non-padding token embeddings of the \texttt{DeepChem/ChemBERTa-77M-MTR} checkpoint, producing 384 features \citep{Ahmad2022}. MHG-GNNs are a GIN-based autoencoder pre-trained on 1.34 million PubChem molecules using $\beta$-VAE loss, producing 1024-dimensional embeddings through iterative message passing and found in the published \texttt{mhggnn\_pretrained\_model\_0724\_2023} checkpoint \citep{kishimoto2023}. We z-score normalized PDV, ChemBERTa and MHG-GNN per feature, using the training molecules' mean and standard deviation; ECFP4, Avalon and Sort \& Slice reach the model as built.

\subsection{Models}
% [notes moved to thesis_notes_moved.md]

% TODO: review, needs manual review
To understand model behavior in the presence of label noise, we explored a range of model families. We fitted tree ensembles (random forests, quantile regression forests, XGBoost, LightGBM and NGBoost), a support vector machine, Gaussian processes and neural networks. Several are variants of one another, differing in the probabilistic machinery they add. All hyperparameters for the models below are included in Additional file~\ref{app:noise:sec:hyperparameters_default}.

While we used 100 trees for RF, we observed a noisy quantile estimate with 100 trees for QRF, prompting us to use 300 instead. To observe the effect of the quantile mechanism, we then fit a second RF at 300 trees so the two models were idential in all other settings.

NGBoost relies on scoring off validation data, which differs by dataset. On QM9 the validation split is already held out of every fit. However, on the three assay datasets the validation set is a scaffold-grouped fifth of the rows NGBoost is handed, which already exclude the validation molecules.

We used a radial basis function (RBF) with $C = 1.0$ and $\gamma = \texttt{scale}$ as the kernel, as it accomodates continuous vectors.

We fitted three variants of GPs, each with a constant mean and a scaled kernel. The first uses an RBF kernel and runs on all six representations. The second uses the Tanimoto kernel \citep{Ralaivola2005, moss2020} and was run on ECFP4 alone. The third is a heteroscedastic variant in which the RBF process is given a second network that predicts the observation noise of each molecule. As such, we used a standard GP and only fit on at most $5{,}000$ training molecules, drawn by a seeded sample. Due to the size of the datasets, this cap only affected QM9. We implemented these GP models in \texttt{GPyTorch} \citep{gpytorch}, with the Tanimoto kernel taken from the \texttt{Gauche} framework \citep{gauche}.

While preliminary tests indicated that the Tanimoto kernel provides GP with substantial benefits over other kernels, for SVM, the RBF kernel was sufficient for all types of data.

% TODO: review, needs manual review
We tested two deterministic feed-forward architectures, both fully connected and both implemented in PyTorch \citep{pytorchGeometric}: (i) NN-$\alpha$ with two hidden layers of sizes 128 and 64, and dropout after each hidden layer and (ii) NN-$\beta$ with two hidden layers, each of size 128, with dropout applied before the output layer. Each model was trained using the Adam optimizer with early stopping on validation loss. We used Bayesian transformations to convert NNs into Bayesian neural networks (BNNs), with ReLU activations and dropout regularization ($p=0.2$). Additional file~\ref{app:noise:sec:hyperparameters_default} gives their full widths, dropout fractions, learning rates and batch size.

To get these networks producing uncertainty estimates, we applied four probabilistic transformations to each architecture, producing eight probabilistic networks. The first replaces every linear layer with a Bayesian layer carrying a Gaussian prior $\mathcal{N}(0, 0.1^2)$ on its weights, producing BNN-$\alpha$ and BNN-$\beta$. They are trained on an evidence lower bound whose  Kullback–Leibler (KL) divergence term is divided by the number of training molecules. While this is a straightforward approach, it is computationally expensive. We explored alternative techniques like our second variant, where we add an output head to the standard NN and fit the negative log likelihood of a Gaussian with that variance \citep{kendall2017}. The third, Variational Bayesian Last Layer (VBLL), replaces every linear layer with a layer that produces mean-field variational posteriors $q(\mathbf{W}) = \mathcal{N}(\boldsymbol{\mu}_W, \text{diag}(\boldsymbol{\sigma}_W^2))$ over the last-layer weights and trains by maximizing the evidence lower bound. This can be thought of as minimizing the reconstruction loss plus KL divergence $D_{\text{KL}}(q(\mathbf{W}) \| p(\mathbf{W}))$ from a standard normal prior, scaled by $1/N$ \citep{Harrison2024}. Because only its output layer keeps a learned observation-noise parameter, VBLL networks report only epistemic uncertainty. The fourth makes observation noise a function of the input. All eight probabilistic networks estimate their predictive distributions from 100 stochastic forward passes. Each network and its Bayesian counterpart were tuned to the same set of parameters so that NN-$\alpha$ against BNN-$\alpha$ and NN-$\beta$ against BNN-$\beta$ differ by the Bayesian layers alone.

\subsection{Noise conditions and levels}
% [notes moved to thesis_notes_moved.md]
% TODO: review, needs manual review
The seven noise conditions are defined in Section~\ref{sec:noise_injection}.

We ran experiments across seven noise levels, $0$, $0.2$, $0.3$, $0.5$, $0.75$, $1.0$ and $1.5$. We added the noise to the raw label and standardized the target afterward, using the clean training mean and spread. For censoring, the level is the fraction of labels clipped, at $0$, $0.1$, $0.2$, $0.25$, $0.3$, $0.4$ and $0.5$ \citep{Svensson2025}.

We ran Gaussian, grouped-wider and grouped-shifted noise on every model and representation pairing. The
remaining four conditions ran on selected pairs. 


\subsection{The QM9 precompute pipeline}
% [notes moved to thesis_notes_moved.md]

% TODO: review, needs manual review
QM9 \citep{Ramakrishnan2014} runs through a separate pipeline, loaded via PyTorch Geometric. Python
loads the data, splits it, computes every representation, and trains and scores the models. A Rust
program sits between featurization and training. It injects the noise, standardizes the labels and
writes the files the models are trained from.

% TODO: review, needs manual review
Each replicate draws a fresh scaffold split, as described under Datasets. Each acyclic molecule's
group is keyed on its canonical SMILES, and a group is never divided between two splits. The same
grouping defines the scaffold families used by the grouped noise conditions.

% TODO: review, needs manual review
Python computes the six representations described above. Binary fingerprints are bit-packed, Sort \&
Slice counts are stored as 16-bit integers, and the descriptors and learned embeddings are stored
as 32-bit floats. A molecule whose ECFP4 or Avalon fingerprint is entirely zero stops the run, so it
cannot enter training as if it had been featurized. One-hot encoded SMILES can still be built but
are not used in this thesis.

% TODO: review, needs manual review
Traditional approaches load entire datasets into memory from CSV or Parquet files, creating RAM bottlenecks at scale and redundant deserialization overhead across experimental replicates. Our pipeline instead serializes preprocessed data into custom memory-mapped binary files (\texttt{.mmap}) that enable on-demand access with constant low memory footprint regardless of dataset size (Figure~\ref{fig:fig_mmap}). Python writes one file per split, the Rust program reads each one and writes a new file with the processed labels in its place, and Python then reads the new files back for training.

% TODO: review, needs manual review
Each mmap file stores a sequence of fixed-format entries. Entry structure includes: length-prefixed UTF-8 SMILES strings (isomeric and canonical), the raw target value (float32, 4 bytes), the requested molecular representations as fixed-width blocks (Sort \& Slice: uint16[1024] = 2,048 bytes, PDV: float32[200] = 800 bytes, ChemBERTa: float32[384] = 1,536 bytes, MHG-GNN: float32[1024] = 4,096 bytes, Avalon: 256 bytes, ECFP4: 256 bytes), and the processed target value incorporating noise and normalization (float32, 4 bytes). Only the requested representations are written, so the byte offset of every field depends on the list of representations. This format enables zero-copy memory mapping—Python and Rust access the same disk-backed memory pages without deserialization, maintaining constant RAM usage regardless of dataset size.

% TODO: review, needs manual review
The Rust program builds a noise plan for each split separately, in that split's own record order.
The training plan defines the condition. The dose, the censoring limit and the affected scaffold
families are all computed from the clean training labels. Validation labels receive an independent
draw, from a seed derived from the training seed, at the same dose in label units, the same limit
and the same affected families. The program checks that validation received the same amount of
noise as training. Test labels are never noised. Every split still records the per-molecule scale
its region would have received, so held-out molecules can be scored against the training pattern.

% TODO: review, needs manual review
For every molecule in every split, the Rust program writes a provenance row: its canonical SMILES,
clean label, injected noise, noise scale, the pattern its region was exposed to, noisy label and the
standardized value written. It also writes a run-level manifest of the condition's parameters and
the delivered dose. The record at each row index must hold the same molecule on the way out as when
the noise was drawn, or the run stops. A self-test mode runs the same dose and targeting checks on a
label file without the rest of the pipeline.

% TODO: review, needs manual review
\paragraph{Models.} The pipeline trains the models described under Models above.
Hyperparameters are either package defaults, values tuned by Optuna \citep{akiba2019optuna}, or
supplied from a JSON file. For models that return uncertainty, the pipeline can also score the
training molecules out of fold, so that each training molecule's uncertainty comes from a model that
never saw its corrupted label.

\begin{figure}[htbp]
    \centering
    \includegraphics[width=0.95\textwidth]{figure_mmap_structure.png}
    \caption[Memory-mapped loading compared with in-memory loading.]{Memory-mapped file architecture contrasted with traditional in-memory loading. Top: CSV/Parquet approaches load entire datasets into RAM, creating bottlenecks at scale. Bottom: Memory-mapped binary files enable on-demand line-by-line access with constant low RAM usage. Right: Entry-level structure showing byte-aligned fields with fixed-size representations (target values, fingerprints, domain labels) and length-prefixed variable-length fields (SMILES strings, one-hot encodings).}
    \label{fig:fig_mmap}
\end{figure}


\subsection{Uncertainty}
% [notes moved to thesis_notes_moved.md]

Up to this point we have been concerned with how noise affects the predictive accuracy of QSAR models. However, label noise should also affect the uncertainty these models report. Some models return only a point prediction, while others also return an estimate of uncertainty, either as a predicted distribution or as the spread across an ensemble. Furthermore, some report a single value of uncertainty while others report one that can be decomposed into an aleatoric and an epistemic component \citep{kendall2017}. The aleatoric component describes underlying noise the model associates with the data, while the epistemic component represents the uncertainty within the model's fit. We sought to determine if, as expected, the aleatoric component rises with injected label noise while the epistemic component does not.

For the BNN and VBLL transformations, we took the two components over $T = 100$ stochastic forward
passes \citep{kendall2017}:
\begin{equation}
\hat{u}^2_{\text{ale}}(x) = \frac{1}{T}\sum_{t=1}^{T}\hat{v}_t(x),
\qquad
\hat{u}^2_{\text{epi}}(x) = \frac{1}{T}\sum_{t=1}^{T}\big(\hat{\mu}_t(x) - \bar{\mu}(x)\big)^2 ,
\label{eq:sampling_split}
\end{equation}
where $\hat{\mu}_t(x)$ and $\hat{v}_t(x)$ are the mean and the variance returned for molecule $x$ on
pass $t$, and $\bar{\mu}(x)$ is the mean over passes. Here $\hat{v}_t(x)$ is the predicted observation
noise variance. It varies per molecule in the variance-head and heteroscedastic variants, and is one
number for the whole fit in the VBLL transformation. A full-BNN without a variance head returns no
variance, so its uncertainty is epistemic alone.

The other models follow the same split. For GPs, the epistemic component is the posterior variance of the latent function and the aleatoric component is the learned observation noise \citep{Rasmussen2005}. All variants of the GP produce one number for the whole fit except heteroscedastic. For QRF models, each tree holds multiple training molecules. So for a given molecule, the aleatoric component is the variance the labels the molecule's given leaf. The total aleatoric uncertainty are those variances averaged over the trees. The epistemic component is the variance of the leaf means across the trees. NGBoost predicts one variance for each molecule from a single fit, which is aleatoric alone \citep{Duan2020}. In every model, the total uncertainty is the sum of the the aleatoric and epistemic components.

No held-out label carries injected noise, so every comparison against the injected noise was made on
training molecules. Each training set was divided into five folds grouped on Murcko scaffolds, and each
molecule was scored by a refit that excluded its own fold. For NGBoost on QM9, three of the five folds
were scored to limit its cost. No uncertainty was calibrated after fitting.


\subsection{Evaluation}
% [notes moved to thesis_notes_moved.md]
We compute statistics on the held-out molecules, for QM9 the test split and for three assay datasets the held-out fold of the scaffold cross-validation. Different combinations of dataset, model, representation, noise condition, and noise level are examined in each scenario. We used the coefficient of determination ($R^2$) throughout. For molecule $i$, let $y_i$ be the clean label, $\hat{y}_i$ the prediction, $\epsilon_i$ the noise injected into its label and $u_i$ the predicted uncertainty.

To evaluate the effect of label noise, we tracked predictive performance retention under increasing artificial noise. For each QSAR pairing we recorded $R^2(\tau)$ at each noise level $\tau$, and normalized it to obtain the retention ratio $R^2(\tau)/R^2(0)$. To quantify robustness, we used the normalized area under this retention curve, \aucnorm (Equation~\ref{eq:auc_norm}), evaluated using the trapezoidal rule over the seven noise levels. While we tested other metrics of robustness, we found \aucnorm to be less directly correlated to clean-label $R^2$ predictions. For censoring, we sweep over a set of fractions to clip the labels prior to normalization. 

An \aucnorm value near 1 indicates that the QSAR model in question has retained predictive performance with artificial noise, meaning it is more robust to noise. Since we normalize the curve by its own clean-label predictive performance, a model with a low clean $R^2$ has less to lose, so \aucnorm is reported beside clean $R^2$. We excluded replicates with a clean $R^2 < 0.3$ from the robustness analysis (Additional file~\ref{app:noise:sec:excluded_configurations}), as performance retention ratios become unstable when the clean-label denominator approaches zero. On QM9 one \aucnorm is computed per replicate, and we report the median over the ten replicates. On the three assay datasets one value is computed per scaffold fold, and the median over the five folds is reported. 

To compare the relative effect of the thirteen base models across datasets and noise conditions, we compared their rankings using Spearman's rank correlation coefficient ($\rho$) between clean $R^2$ and \aucnorm. We compared probabilistic models  with its deterministic counterpart using a two-sided Wilcoxon signed-rank test paired on the replicate, at $\alpha = 0.05$, separately for each representation and noise condition. This test was run only on QM9, as with five scaffold folds the
smallest attainable $p$ is 0.0625.

We conducted a separate two-way analysis of variance (ANOVA) for each noise condition to identify the
relative contributions of molecular representation and model architecture to both predictive
performance and noise robustness. For a given metric $y$, either $R^2$ at a noise level of 1.0 or
\aucnorm:
\begin{equation}
y_{ijr} = \mu + \alpha_i + \beta_j + (\alpha\beta)_{ij} + e_{ijr},
\end{equation}
where $\alpha_i$ represents the effect of model architecture $i$, $\beta_j$ represents the effect of
molecular representation $j$, $(\alpha\beta)_{ij}$ is the interaction term, and $e_{ijr}$ is the
residual for replicate $r$. The proportion of variance explained by each term was calculated as the
effect size
\begin{equation}
\eta^2_{\text{term}} = \frac{SS_{\text{term}}}{SS_{\text{total}}},
\end{equation}
using Type I (sequential) sums of squares, entered in the order model architecture, molecular
representation, interaction. The ten QM9 replicates provide the residual. The uncertainty on each share is the range of $\eta^2$ when the decomposition is repeated with each replicate left out in turn. 

ANOVA relies on the assumption that observations both within and across groups are independence. Across the whole set of models and representations in this study, some pairs produced similar predictions that are not independent. This causes the residual degrees of freedom to increase, while the residual sum of squares barely moves. To handle this issue, we computed pairwise Spearman rank correlations between all models and between all representations across noise strategies. Models whose profiles correlated at $\rho > 0.99$ with another model were excluded, including quantile regression forests (QRF), which were very similar to RF, as well as different BNN variants. 

Across representations, SNS correlated at $\rho > 0.90$ with the ECFP4 fingerprint and was excluded to avoid inflating representation degrees of freedom, as both encode overlapping circular substructure information. We also computed intraclass correlation coefficients ICC(1,1) for all model pairs. ICC(1,1) measures the proportion of total variance attributable to between-subject differences, with values near 1.0 indicating that two models rank configurations almost identically. The full redundancy and ICC tables are reported in Additional files~\ref{app:noise:sec:model_redundancy}--\ref{app:noise:sec:model_icc}.

To assess the interaction between uncertainty estimation and label noise, we posed three questions. First, whether uncertainty rise with label noise, as measured by the mean $u_i$ over held-out molecules at the highest noise level
divided by its value with no added noise, in each replicate or fold. The second pertains to models who report separate aleatoric and epistemic components. We divided the increase of mean epistemic uncertainty between the highest level of noise and clean labels by the mean of the aleatoric component to determine how each responds to label noise. A ratio near 0 means the added noise was attributed to the labels while a ratio of 1 or above means the epistemic component grew as much or more. We also compared the mean aleatoric component with the standard deviation of the injected noise. The third question we had was whether per-sample uncertainty values track injected label noise, as measured by the Spearman correlation between $u_i$ and $|\epsilon_i|$ over out-of-fold training molecules, per fold and per noise level. We also report the Spearman correlation between $u_i$ and the absolute error of the prediction. Finally, we asked whether dividing a molecule's out-of-fold absolute error by $u_i$ ranks the corrupted labels better than the error alone, and report the difference between the two Spearman correlations with $|\epsilon_i|$.


\section{Representation governs clean performance, architecture governs robustness}
% [notes moved to thesis_notes_moved.md]


Considering the two main components of a QSAR model, the molecular representation and the model architecture, we sought to understand how each contributes to predictive performance and noise robustness. ANOVA decomposition across six noise strategies on models to predict HOMO--LUMO gap from QM9 demonstrates that for predictive performance (R$^2$ at $\sigma = 0.3$), the model--representation interaction term is the  largest source of variance across all six strategies, followed by model architecture and representation (Table~\ref{tab:anova_decomposition}; Figure~\ref{fig:anova_decomposition}a). However, for noise robustness as measured by the noise degradation slope (NDS), this ordering shifts. The choice of model architecture is instead the largest source of variance, while the interaction term drops to second and representation third. This observed inversion makes sense because the molecular representation dictates the chemical information the model receives, while the architecture determines how efficiently that information is used. The interaction term reflects the idea that these two are not independent. In the context of label noise robustness, representation has less of an impact since the features are not corrupted by noise. Meanwhile, the model's regularization, ensembling, and uncertainty handling have a stronger impact on how it handles noise. 

The impact varied by type of noise, represented by noise strategy (Figure~\ref{fig:anova_decomposition}b). The strategies which degraded in performance the most, threshold and value-proportional, both corrupted label regions, forcing models to rely on their inductive biases. This translates to the model architecture explaining a significant portion of the variance. On the other hand, outlier noise only affects statistical outliers. For outlier noise, model and interaction effects were both smaller.

\begin{table}[htbp]
\centering
\caption[ANOVA variance decomposition by noise strategy on QM9.]{ANOVA variance decomposition by noise strategy on the QM9 HOMO--LUMO gap. Each cell reports $\eta^2$ (\%), the share of variance explained by that factor. Columns: \textbf{Model} = model architecture effect; \textbf{Rep} = molecular representation effect; \textbf{Inter.}\ = model$\times$representation interaction; \textbf{Resid.}\ = residual unexplained variance. ``Performance'' uses R$^2$ at $\sigma = 0.3$ as the outcome; ``Robustness'' uses the noise degradation slope (NDS). Bold entries highlight the dominant factor in each row.}
\label{tab:anova_decomposition}
\small
\begin{adjustbox}{max width=\textwidth}
\begin{tabular}{lcccccccc}
\toprule
& \multicolumn{4}{c}{\textbf{Performance ($\eta^2$, \%)}} & \multicolumn{4}{c}{\textbf{Robustness ($\eta^2$, \%)}} \\
\cmidrule(lr){2-5} \cmidrule(lr){6-9}
\textbf{Strategy} & \textbf{Model} & \textbf{Rep} & \textbf{Inter.} & \textbf{Resid.} & \textbf{Model} & \textbf{Rep} & \textbf{Inter.} & \textbf{Resid.} \\
\midrule
Gaussian        & 24.6 & 22.8 & \textbf{48.1} &  4.5 & \textbf{48.7} &  8.2 & 11.9 & 31.2 \\
Quantile        & 25.1 & 22.4 & \textbf{46.9} &  5.6 & \textbf{47.9} &  6.8 & 14.9 & 30.4 \\
Threshold       & 25.0 & 21.2 & \textbf{48.0} &  5.8 & \textbf{48.0} & 10.1 & 28.0 & 13.9 \\
Heteroscedastic & 27.9 & 23.5 & \textbf{46.2} &  2.5 & \textbf{37.0} &  4.5 & 17.5 & 41.0 \\
Value-prop.     & 27.8 & 22.0 & \textbf{46.5} &  3.8 & \textbf{52.8} &  6.9 & 22.2 & 18.2 \\
Outlier         & 25.3 & 22.0 & \textbf{49.2} &  3.5 & 12.7 &  3.2 &  4.7 & \textbf{79.3} \\
\bottomrule
\end{tabular}
\end{adjustbox}
\end{table}

\begin{figure}[htbp]
    \centering
    \includegraphics[width=\textwidth]{fig2_anova_decomposition.png}
    \caption[ANOVA variance decomposition of performance and robustness on QM9.]{ANOVA variance decomposition ($\eta^2$, \%) for a) predictive performance (R$^2$ at $\sigma=0.3$) and b) noise robustness, quantified by the noise degradation slope (NDS), by noise strategy. Lower (more negative) NDS indicates faster performance loss under noise.}
    \label{fig:anova_decomposition}
\end{figure}

The interaction between model architecture and representation is seen in Figure~\ref{fig:interaction}. While some model architectures, like SVM and full BNNs maintain consistent noise robustness scores across all representations, others like NN-$\beta$ display greater variation. Model robustness on one representation strongly predicts robustness on another (Spearman $\rho = 0.73$, $p = 0.007$ for ECFP4 vs PDV; 11 models). However, the scatter around the diagonal reflects the interaction term identified by the ANOVA. 

\begin{figure}[htbp]
\centering
\includegraphics[width=\textwidth]{fig_interaction.png}
\caption[Representation--model interaction effects on QM9.]{Representation--model interaction effects on the QM9 HOMO--LUMO gap. a) Noise degradation slope (NDS) heatmap showing model robustness across representations under Gaussian noise; lower (more negative) NDS indicates faster degradation. ``N/A'' indicates configurations excluded due to baseline $R^2 < 0.6$ or kernel incompatibility (GP uses a Tanimoto kernel defined on binary vectors, which is incompatible with PDV). b) NDS on PDV versus ECFP4 fingerprint for each model under Gaussian noise, with Spearman correlation.}
\label{fig:interaction}
\end{figure}

We filtered out model--representation configurations with baseline predictive performance R$^2 < 0.6$, excluding 66 configurations where poor clean-data performance would produce misleadingly shallow degradation slopes (Additional file~\ref{app:noise:sec:excluded_configurations}). We evaluated all model architectures under increasing label noise on the PDV representation across six strategies, considering both predictive performance as measured by $R^2$ as well as the noise degradation slope (NDS). We observe that NDS varies substantially across architectures and noise strategies, as seen in Figure~\ref{fig:global_overview}. 

\begin{figure}[htbp]
\centering
\includegraphics[width=\textwidth]{fig1_global_overview.png}
\caption[Overview of noise robustness on QM9.]{Global overview of noise robustness on the QM9 HOMO--LUMO gap. a) Performance degradation curves (PDV, Gaussian noise) showing R$^2$ versus noise level $\sigma$ for representative models. b) Noise degradation slope (NDS) heatmap across all model architectures and noise strategies on PDV; lower (more negative) NDS indicates faster degradation under noise. The Tanimoto GP is absent because its kernel is incompatible with PDV (see Section~\ref{sec2}); a separate GP with an RBF kernel is included for PDV-specific analyzes.}
\label{fig:global_overview}
\end{figure}

\section{Probabilistic models are no more noise-robust than deterministic ones}
% [notes moved to thesis_notes_moved.md]

To further explore the impact of model architecture on noise robustness, we examined variants of the same core architecture. Specifically, Bayesian transformations were used to explore the impact of the Bayesian framework on noise. By comparing models with the same core architecture, we observe that with full Bayesian transformations, representation explains less than $11\%$ of NDS variance ($p > 0.25$) while variational Bayesian inference on the same NNs is representation-dependent. On MHG-GNN and mol2vec representations, VBLL models were unstable, such that all 24 VBLL $\times$ \{MHG-GNN, mol2vec\} configurations were excluded due to baseline R$^2 < 0.6$. Further work could involve more thoroughly-tuning these models and exploring different core architectures to adapt to these embedding-style inputs. 

In general, the full Bayesian transformations reduced the degradation in performance with Gaussian noise the most, as seen in Figure~\ref{fig:nn_family_comparison}. Both full BNN and VBLL transformations significantly improved robustness for both NN-$\alpha$ and NN-$\beta$ (Table~\ref{tab:wilcoxon_bnn}). Full transformations are more computationally expensive, making VBLL a potentially stronger approach to handling noise with NNs when working with larger data sets or modeling architectures. The choice of transformation varies by representation, as seen in Additional file~\ref{app:noise:sec:bayesian_transformation_tests}, with SMILES receiving the largest benefits from Bayesian transformations. QRF was significantly less robust than RF (Table~\ref{tab:wilcoxon_bnn}), suggesting that quantile regression adds complexity and flexibility which may result in overfitting.

\begin{table}[htbp]
\centering
\caption[Wilcoxon signed-rank tests for the probabilistic transformations on QM9.]{Wilcoxon signed-rank tests for Bayesian and probabilistic transformations on the QM9 HOMO--LUMO gap. \textbf{Base} is the deterministic reference model; \textbf{Variant} is the probabilistic/Bayesian transformation tested against it. \textbf{$\Delta$ NDS} = mean change in noise degradation slope (NDS) from base to variant; positive values mean the variant is more robust (slope closer to zero). \textbf{Sig.}\ marks pairs that are statistically significant at $p < 0.05$ (asterisk).}
\label{tab:wilcoxon_bnn}
\small
\begin{adjustbox}{max width=\textwidth}
\begin{tabular}{llrrl}
\toprule
\textbf{Base} & \textbf{Variant} & \textbf{$\Delta$ NDS} & \textbf{$p$-value} & \textbf{Sig.} \\
\midrule
NN-$\alpha$ & BNN-$\alpha$ & $+0.056$ & $2.9 \times 10^{-10}$ & \textbf{*} \\
NN-$\alpha$ & VBLL-$\alpha$ & $+0.061$ & $1.2 \times 10^{-6}$ & \textbf{*} \\
NN-$\beta$ & BNN-$\beta$ & $+0.096$ & $2.9 \times 10^{-11}$ & \textbf{*} \\
NN-$\beta$ & VBLL-$\beta$ & $\mathbf{+0.124}$ & $1.2 \times 10^{-7}$ & \textbf{*} \\
RF  & QRF  & $\mathbf{-0.022}$ & $1.6 \times 10^{-10}$ & \textbf{*} \\
\bottomrule
\end{tabular}
\end{adjustbox}
\end{table}

\begin{figure}[htbp]
    \centering
    \includegraphics[width=\textwidth]{fig_nn_family_comparison.png}
    \caption[Neural network families under Gaussian noise on QM9.]{Neural network family comparison on the QM9 HOMO--LUMO gap with N=10k: R$^2$ versus $\sigma$ under Gaussian noise on the physicochemical descriptor vector (PDV) representation for a) NN-$\alpha$ family, b) NN-$\beta$ family, and c) RF versus QRF. Heteroscedastic results (not shown) are qualitatively consistent.}
    \label{fig:nn_family_comparison}
\end{figure}

\section{Relative model robustness is largely strategy-independent}
% [notes moved to thesis_notes_moved.md]

NGBoost and SVM showed the smallest degradation in predictive performance across all strategies, while NN-$\beta$ showed the steepest (Table~\ref{tab:nds_ranking}). When sorting models by NDS, their ranks across noise strategies were consistent: Kendall's $W = 0.92$ ($p = 2.7 \times 10^{-8}$; 11 models, 6 strategies), indicating that a model's relative robustness is largely strategy-independent. The same ranking pattern holds on the ECFP4 fingerprint (Additional file~\ref{app:noise:sec:ecfp4_overview}). Robustness is also largely decoupled from baseline predictive performance: across models on PDV under Gaussian noise, NDS clusters near $-0.38$ regardless of baseline R$^2$ (Additional file~\ref{app:noise:sec:clean_r2_vs_robustness}), so a high-performing model is not automatically a more robust one.

\begin{table}[htbp]
\centering
\caption[Noise degradation slope by model on PDV, QM9.]{Noise degradation slope (NDS) by model on PDV (QM9 HOMO--LUMO gap), ranked by mean across six strategies. Lower $|$NDS$|$ (slope closer to zero) indicates greater robustness. Strategy abbreviations: \textbf{Gauss.}\ = Gaussian, \textbf{Quant.}\ = Quantile, \textbf{Thresh.}\ = Threshold, \textbf{Hetero.}\ = Heteroscedastic, \textbf{Val.-P.}\ = Value-Proportional. Bold rows mark the two most robust models; bold cell entries mark each strategy's most robust configuration.}
\label{tab:nds_ranking}
\small
\begin{adjustbox}{max width=\textwidth}
\begin{tabular}{lccccccc}
\toprule
\textbf{Model} & \textbf{Gauss.} & \textbf{Outlier} & \textbf{Quant.} & \textbf{Thresh.} & \textbf{Hetero.} & \textbf{Val.-P.} & \textbf{Mean} \\
\midrule
\textbf{NGBoost}    & $\mathbf{-0.33}$ & $-0.12$ & $\mathbf{-0.29}$ & $\mathbf{-0.65}$ & $\mathbf{-0.18}$ & $\mathbf{-0.58}$ & $\mathbf{-0.36}$ \\
\textbf{SVM}        & $-0.36$ & $\mathbf{-0.11}$ & $-0.30$ & $-0.71$ & $-0.20$ & $-0.63$ & $-0.38$ \\
RF                  & $-0.36$ & $-0.12$ & $-0.30$ & $-0.71$ & $-0.20$ & $-0.63$ & $-0.39$ \\
LightGBM            & $-0.36$ & $-0.13$ & $-0.32$ & $-0.72$ & $-0.20$ & $-0.64$ & $-0.40$ \\
VBLL-$\alpha$       & $-0.38$ & $-0.16$ & $-0.35$ & $-0.68$ & $-0.22$ & $-0.60$ & $-0.40$ \\
XGBoost             & $-0.37$ & $-0.13$ & $-0.33$ & $-0.72$ & $-0.20$ & $-0.65$ & $-0.40$ \\
VBLL-$\beta$        & $-0.39$ & $-0.14$ & $-0.34$ & $-0.70$ & $-0.23$ & $-0.64$ & $-0.41$ \\
BNN-$\alpha$        & $-0.38$ & $-0.15$ & $-0.35$ & $-0.72$ & $-0.21$ & $-0.64$ & $-0.41$ \\
BNN-$\beta$         & $-0.38$ & $-0.15$ & $-0.36$ & $-0.72$ & $-0.21$ & $-0.65$ & $-0.41$ \\
NN-$\alpha$         & $-0.41$ & $-0.16$ & $-0.37$ & $-0.77$ & $-0.23$ & $-0.68$ & $-0.44$ \\
NN-$\beta$          & $-0.49$ & $-0.17$ & $-0.42$ & $-0.88$ & $-0.26$ & $-0.79$ & $-0.50$ \\
\bottomrule
\end{tabular}
\end{adjustbox}
\end{table}

The strategies differed not only in predictive performance severity, but also in their ability to discriminate between architectures. The spread of NDS across models differed by noise strategy, with outlier noise having a spread of 0.06 and threshold noise having 0.23, presumably exposing architectural differences that are masked by other types of noise. As seen in Table~\ref{tab:nds_ranking}, NGBoost was the most resistant to noise across strategies. With respect to the Variational Bayesian Last Layer (VBLL) models, VBLL-$\alpha$ showed a similar pattern, outperforming all tree-based models except NGBoost under threshold and value-proportional noise despite ranking below them under Gaussian noise. This suggests that its learned noise variance is most beneficial when label corruption is systematic rather than uniform. This observation with VBLL is not alone; predictive performance on clean data is not well-correlated with noise robustness. In fact, some of the models with the lowest NDS across strategies such as NGBoost and SVM did not perform particularly well on clean data. These models may be more prone to picking up on noise rather than underlying signals in the data. 

% TODO: review, needs manual review
As established by the ANOVA, model architecture and the model--representation interaction term dominate NDS variance. A handful of models stood out as being noise-robust regardless of representation, namely SVM and full BNNs. Both methods rely on inductive biases, SVM on margin maximization and BNN on weight priors. However, other models like RF and NN-$\beta$ show the opposite pattern. These models can be robust to noise, but only when paired with particular representations. The pairings differ by model. 

% TODO: review, needs manual review
While representation may not be a dominant determinant in NDS variance, the model--representation interaction term does have an impact. For SMILES, model choice explains over 91\% of robustness variance, while on PDV, model choice explains 72\%. Across all six noise strategies, PDV configurations produced the shallowest mean NDSs across representations, with its best performance coming with models with strong inductive biases like SVM and full BNN. This suggests that PDVs are preferable when dealing with noisy data. It appears that PDV's compact feature space complements margin-based and Bayesian regularization. Embeddings struggled, particularly when paired with NNs. These patterns are aligned with the importance of the model--representation interaction term, indicating that representation choice modulates noise robustness within a given architecture. 

Taking a closer look at model choice for noise robustness, we ranked all model--representation configurations on QM9 by their Gaussian NDS. Tree-based ensembles dominate the top 10: NGBoost appears five times, XGBoost and RF twice each, and LightGBM once. No neural-network configuration appears the top 10. Rankings are consistent across strategies and the top QM9 configurations maintain their advantage on ADME datasets, though SVM configurations rank higher on these ADME datasets than NGBoost.

\section{The pattern holds across prediction targets}
% [notes moved to thesis_notes_moved.md]

\begin{figure*}[htbp]
\centering
\includegraphics[width=\textwidth]{fig_validation_overview.png}
\caption[Noise degradation slope on LogD, Caco-2 and hERG $K_i$.]{Noise degradation slope (NDS) heatmaps for three external validation datasets: a) LogD, b) Caco-2 Efflux, and c) hERG K$_i$. Lower (more negative) NDS indicates faster degradation under noise. Black cells indicate baseline R$^2 < 0.3$ for that (model, strategy) pair; ``N/A'' indicates filtered extreme values ($|$NDS$| > 2$). Models whose baseline R$^2$ fell below 0.3 across all strategies are omitted entirely, accounting for the smaller set of rows on hERG K$_i$.}
\label{fig:validation_overview}
\end{figure*}

In addition to working with ``clean'' quantum molecular properties, we tested our findings on experimentally-obtained data, replicating our process of adding artificial noise. Unlike QM9, where $\sigma = 0$ corresponds to a noise-free baseline, the experimental endpoints carry their own measurement noise in both training and test labels; the injected $\sigma$ therefore stacks on top of an unknown noise floor. Per-dataset NDS heatmaps across all model architectures and noise strategies are shown in Figure~\ref{fig:validation_overview}. The ANOVA on these additional datasets confirms that model architecture dominates robustness variance on all three datasets, as seen in Additional file~\ref{app:noise:sec:assay_variance_decomposition}, and that the trends with respect to noise robustness generalize. NGBoost ranks first on both QM9 and external datasets, SVM also transfer quite well. Some models' predictive capabilities are limited on these external datasets, likely due to a combination of smaller training sets, narrower chemical coverage, and the unknown experimental noise in the labels. XGBoost and BNN variants suffer the most here, as seen in Figure~\ref{fig:validation_combined}a.

\begin{figure*}[htbp]
\centering
\includegraphics[width=\textwidth]{fig_validation_combined.png}
\caption[Noise robustness on QM9 and three assay datasets.]{Comparison of noise robustness on QM9 (HOMO--LUMO gap, $N = 10{,}000$) and three external validation datasets: LogD ($N = 5{,}039$), Caco-2 Efflux ($N = 2{,}161$), and hERG K$_i$ ($N = 1{,}482$). a) Model robustness measured by the noise degradation slope (NDS) across all four datasets; lower (more negative) NDS indicates faster degradation under noise. b) QM9 vs external robustness: models with inherent robustness mechanisms (NGBoost, SVM) maintain their relative standing, while ensemble-dependent models (XGBoost) degrade on the external datasets.}
\label{fig:validation_combined}
\end{figure*}

QRF was consistently less robust than RF on every external data set (Additional file~\ref{app:noise:sec:rf_vs_qrf_assay}). As seen with QM9, choice of molecular representation has a minimal effect on noise robustness across all three external datasets, reinforcing the finding that model architecture, not representation, is the main driver of noise robustness. 

\section{Uncertainty under label noise}\label{sec:uq}
% [notes moved to thesis_notes_moved.md]
% TODO: Hirschfeld2020, Scalia2020 and Norinder2014 are in refs.bib but not cited anywhere in the thesis. Cite them in this section or drop them.

The sections above showcased the importance of model selection when working with noisy labels. We take this one step further and ask, for the models which produce uncertainty estimates, how those estimates are affected by noise and if they can be used to track noise. 

\subsection{Models with an explicit noise term track prediction error under noise}
% [notes moved to thesis_notes_moved.md]

\citet{Kolmar2021} found that the mean predicted uncertainty derived from the GPs increases with the amount of label noise in the training data. We go beyond the population level and instead ask if a model's per-sample uncertainty tracks label noise. One key sign that a model is able to handle noise is its ability to track it, represented by the per-sample Spearman correlation between predicted uncertainty and noise magnitude (Table~\ref{tab:top_unc_noise}). While GPs and NGBoost gave the strongest correlations, BNNs were moderate. Learned embeddings (MHG-GNN and mol2vec) struggled, so much so that their correlations dropped to near zero or below on every model tested. 

\begin{table}[htbp]
\centering
\caption[Strongest and weakest combinations for uncertainty--noise correlation on QM9.]{Strongest and weakest model--representation combinations for uncertainty--noise correlation on the QM9 HOMO--LUMO gap (Gaussian noise strategy). Columns: \textbf{Unc-Noise $\rho$} = Spearman correlation between predicted uncertainty and injected noise magnitude $|\epsilon|$ (higher = better noise detection); \textbf{Unc-Error $\rho$} = Spearman correlation between predicted uncertainty and absolute prediction error; \textbf{ECE} = expected calibration error (lower = better-calibrated); \textbf{Cov.\ $1\sigma$} / \textbf{$2\sigma$} = empirical coverage at the $1\sigma$ (68\% target) and $2\sigma$ (95\% target) intervals.}
\label{tab:top_unc_noise}
\small
\begin{adjustbox}{max width=\textwidth}
\begin{tabular}{llrrrrr}
\toprule
\textbf{Model} & \textbf{Rep.} & \textbf{Unc-Noise $\rho$} & \textbf{Unc-Error $\rho$} & \textbf{ECE} & \textbf{Cov. $1\sigma$} & \textbf{Cov. $2\sigma$} \\
\midrule
\textbf{GP}    & \textbf{SNS}    & \textbf{0.56} & 0.23 & 0.14 & 71.2\% & 94.6\% \\
GP             & Morgan          & 0.53 & 0.22 & 0.14 & 71.4\% & 95.4\% \\
\textbf{NGBoost} & \textbf{SNS}  & \textbf{0.47} & 0.16 & 0.13 & 69.3\% & 95.4\% \\
\textbf{NGBoost} & \textbf{PDV}  & \textbf{0.47} & \textbf{0.29} & \textbf{0.11} & 69.0\% & 95.2\% \\
NGBoost        & PDV (binary)    & 0.44 & 0.25 & 0.13 & 69.8\% & 95.4\% \\
NGBoost        & Morgan          & 0.44 & 0.20 & 0.14 & 69.8\% & 95.3\% \\
BNN-$\alpha$   & SNS             & 0.42 & 0.21 & 0.17 & 73.4\% & 95.1\% \\
BNN-$\beta$    & SNS             & 0.40 & 0.21 & 0.18 & 73.5\% & 95.1\% \\
BNN-$\beta$    & Morgan          & 0.38 & 0.20 & 0.18 & 72.1\% & 95.0\% \\
BNN-$\beta$    & PDV             & 0.38 & 0.24 & 0.19 & 72.1\% & 94.0\% \\
\midrule
\multicolumn{7}{l}{\textit{Contrast: lowest correlations (graph/embedding representations)}} \\
\midrule
BNN-$\beta$ & mol2vec & $-0.05$ & 0.02 & 1.33 & 50.6\% & 76.4\% \\
VBLL-$\alpha$ & MHG-GNN & $-0.05$ & 0.09 & 13.41 & 26.6\% & 38.6\% \\
VBLL-$\beta$ & MHG-GNN & $-0.03$ & 0.10 & 15.34 & 31.1\% & 47.1\% \\
VBLL-$\beta$ & mol2vec & $-0.03$ & 0.15 & 23.86 & 44.8\% & 58.0\% \\
\bottomrule
\end{tabular}
\end{adjustbox}
\end{table}

\subsection{NGBoost and the GP stay near target coverage while Bayesian variants miscalibrate}
% [notes moved to thesis_notes_moved.md]

Samples containing label noise tend to produce large residuals in training, and every model has its own way of handling them. NGBoost contains predicted scales that increase to absorb residual \citep{Duan2020}. The GP contains has a single global observation noise term. It does not absorb individual residuals, so its per-sample uncertainty derives primarily from the posterior variance rather than label noise  \citep{Rasmussen2005, Obrezanova2007}. BNNs treat these large residuals differently; they broaden their weight posteriors, increasing predicted uncertainty without explicitly modeling observation noise \citep{gal2016, kendall2017}. VBLL adds a learned noise variance to the BNN loss \citep{Harrison2024}, but this value is a global scalar, independent of the input. As seen in Table~\ref{tab:top_unc_noise} and Additional file~\ref{app:noise:sec:uncertainty_metrics_ecfp4}, both BNNs and VBLLs achieve moderate correlations which improve when paired with fingerprint or descriptor-based representations. Among models that learn separate distributional parameters during training, GP and NGBoost produced the strongest correlations between uncertainty and label noise, though QRF, which separates parameters using quantile distributions, produced both poor label predictions and uncertainty-noise correlations. At the population level, mean predicted uncertainty increases with artificial noise. Although we expect to see aleatoric uncertainty increase with injected noise, for VBLL both the aleatoric and epistemic components increased (Figure~\ref{fig:uncertainty_combined}b).

\begin{figure}[htbp]
\centering
\includegraphics[width=\textwidth]{fig_uncertainty_combined.png}
\caption[Uncertainty estimates under label noise on QM9.]{Uncertainty estimates under label noise on the QM9 HOMO--LUMO gap. a) Mean predicted uncertainty versus $\sigma$ for probabilistic models (PDV, Gaussian strategy). b) Aleatoric and epistemic uncertainty components versus $\sigma$ for VBLL models.}
\label{fig:uncertainty_combined}
\end{figure}

\subsection{Uncertainty attribution is itself representation-dependent}
% ---- NOTES: Uncertainty attribution is itself representation-dependent ----
% Takeaway: the variance-head BNNs put added noise into the aleatoric part, as expected. QRF puts much of it into the epistemic part, and the representation decides how much.
% Takeaway: the aleatoric part shows that noise rose, not how much. Under censoring, uncertainty shrinks instead of rising.
% Paragraph 1: the question, and which models can answer it.
% - Mean uncertainty rises with noise, but that does not say what the model attributes the rise to.
% - Expected behavior: the aleatoric part rises with injected label noise; the epistemic part stays roughly flat.
% - Three model kinds give a per-molecule split: the full BNNs with a variance head, the heteroscedastic GP, and QRF \citep{kendall2017, Rasmussen2005, Meinshausen2006}.
% - The other GPs are excluded: their observation noise is one number for the whole fit, so it cannot follow individual molecules \citep{Rasmussen2005}.
% - Measure: the rise in mean epistemic uncertainty divided by the rise in mean aleatoric uncertainty, from no added noise to noise level 1.5, over held-out molecules.
% - A ratio near 0 means the added noise went to the aleatoric part. A ratio of 1 or more means the epistemic part grew as much or more.
% Paragraph 2: the variance-head BNNs split as expected; QRF does not.
% - Full BNNs with a variance head: the ratio is at most 0.43 in every model, dataset, representation and noise condition.
% - QRF: both parts rise with noise, and on some representations the epistemic part rises faster.
% - This holds although QRF is among the more noise-robust models: on QM9 under Gaussian noise it ranks 4th of the 19 models and variants on AUC_norm on four of the six representations.
% - Do not call QRF among the most accurate. On the same data its clean R^2 ranks between 6th and 13th of the 19, depending on the representation.
% Paragraph 3: for QRF the representation decides whether the split works.
% - Contrast: representation explained little of the variance in robustness, but here it decides the split.
% - QRF ratio on held-out molecules, range over datasets and every noise condition except censoring: 0.63 to 1.09 on the three fingerprints (ECFP4, Avalon, Sort & Slice).
% - Same range on PDV: 1.03 to 1.42. On ChemBERTa and MHG-GNN: 1.28 to 2.05.
% - Use the held-out range for PDV (1.03 to 1.42), not the paper's 1.01 to 1.45, which pools held-out and out-of-fold training molecules.
% - Under Gaussian noise alone, held-out molecules: ECFP4 0.65-0.75, Sort & Slice 0.69-0.76, Avalon 0.81-0.89, PDV 1.17-1.32, MHG-GNN 1.54-1.72, ChemBERTa 1.60-1.88.
% - So on the learned embeddings QRF puts more of the added noise into its epistemic part than into its aleatoric part.
% Paragraph 4: grouped-shifted noise moves the heteroscedastic GP, most on fingerprints.
% - Heteroscedastic GP, Gaussian noise, held-out molecules: ratio between -0.05 and 0.36 on every representation.
% - Heteroscedastic GP, grouped-shifted, range over datasets: ECFP4 0.73-1.00, Sort & Slice 0.44-0.61, Avalon 0.32-0.46.
% - Same model and condition: PDV 0.13-0.22, ChemBERTa 0.14-0.23, MHG-GNN 0.21-0.25.
% - Laplace, Student-t and outlier noise give splits like those under Gaussian noise.
% - Grouped-shifted raises the epistemic share of the heteroscedastic GP and of QRF for every dataset and representation.
% - Possible reading, stated as a reading: a shared offset per scaffold family looks to the model like structure in the data, not label noise.
% Paragraph 5: the aleatoric part shows that noise rose, not how much.
% - This comparison is against the injected noise, so it uses out-of-fold training molecules. Held-out labels carry no injected noise.
% - Slope of mean aleatoric uncertainty against the standard deviation of the added noise, both in the label's units, Gaussian noise, median over folds.
% - On ECFP4 and PDV: 0.36 (QRF, PDV, Caco-2) to 1.09 (a full BNN with a variance head, PDV, hERG K_i).
% - Change: the paper's text says 0.36 to 0.99. The upper end from the fitted slopes is 1.09; re-check the paper's source.
% - Including ChemBERTa, the range widens to 0.28 (QRF, Caco-2) to 1.62 (a full BNN with a variance head, hERG K_i).
% - So it can compare noise between datasets or noise levels. It does not estimate a dataset's own measurement error.
% Paragraph 6: censoring shrinks the uncertainty instead of raising it.
% - Clipping labels at a limit narrows their spread.
% - NGBoost and QRF shrink their total uncertainty with that spread on every dataset; the GPs do the same on hERG K_i and Caco-2.
% - None of these models treats a clipped label as a noisy one, so they grow more confident while losing information.
% Table: ratio of epistemic rise to aleatoric rise. One row per model (the two full BNNs with a variance head, heteroscedastic GP, QRF) and noise condition (Gaussian, grouped-shifted).
% - One column per representation (ECFP4, Avalon, Sort & Slice, PDV, ChemBERTa, MHG-GNN). Each entry is the range over datasets, held-out molecules.
% Figure: aleatoric and epistemic parts against noise level, one panel per model with a per-molecule split, QM9 HOMO--LUMO gap, ECFP4, Gaussian noise.
% Figure: aleatoric part against the standard deviation of the added noise, per model, ECFP4 and PDV, out-of-fold training molecules, with a dashed line where reported noise equals added noise.
% ---- END NOTES ----

% CODE INPUT
Mean predicted uncertainty rises with label noise, but that alone does not say what the model attributes the rise to. We expected the aleatoric part to rise with injected label noise, and the epistemic part to stay roughly flat. Three kinds of model split their uncertainty per molecule: the full BNNs with a variance head, the heteroscedastic GP, and QRF \citep{kendall2017, Rasmussen2005, Meinshausen2006}. We excluded the other GPs. Their observation noise is one number for the whole fit, so it cannot follow individual molecules \citep{Rasmussen2005}. For each model we took the rise in mean epistemic uncertainty from no added noise to noise level 1.5, over held-out molecules. We divided it by the rise in mean aleatoric uncertainty over the same range. A ratio near 0 means the added noise went to the aleatoric part. A ratio of 1 or more means the epistemic part grew as much or more (\textbf{[FIGURE NEEDED: table of the ratio of epistemic rise to aleatoric rise, one row per model (the two full BNNs with a variance head, heteroscedastic GP, QRF) and noise condition (Gaussian, grouped-shifted), one column per representation, each entry the range over datasets on held-out molecules]}).

% CODE INPUT
The full BNNs with a variance head split the added noise as expected. Their ratio is at most 0.43 in every model, dataset, representation and noise condition. QRF does not. Both of its parts rise with noise, and on some representations the epistemic part rises faster (\textbf{[FIGURE NEEDED: aleatoric and epistemic parts against noise level, one panel per model with a per-molecule split, QM9 HOMO--LUMO gap, ECFP4, Gaussian noise]}). This holds although QRF is among the more noise-robust models. On QM9 under Gaussian noise, it ranks 4th of the 19 models and variants on \aucnorm on four of the six representations. Its clean $R^2$ on the same data ranks between 6th and 13th of the 19, depending on the representation.

% CODE INPUT
Representation explained little of the variance in noise robustness, but for QRF it decides how the added noise is split. Over datasets and every noise condition except censoring, the QRF ratio on held-out molecules ranges from 0.63 to 1.09 on the three fingerprints, ECFP4, Avalon and Sort \& Slice. On PDV the same range is 1.03 to 1.42. On ChemBERTa and MHG-GNN it is 1.28 to 2.05. Under Gaussian noise alone, the held-out ranges are 0.65--0.75 on ECFP4, 0.69--0.76 on Sort \& Slice and 0.81--0.89 on Avalon. They are 1.17--1.32 on PDV, 1.54--1.72 on MHG-GNN and 1.60--1.88 on ChemBERTa. On the learned embeddings, QRF therefore puts more of the added noise into its epistemic part than into its aleatoric part.

% CODE INPUT
The type of noise also moves the split, most clearly for the heteroscedastic GP. Under Gaussian noise its ratio on held-out molecules lies between $-0.05$ and 0.36 on every representation. Under grouped-shifted noise the range over datasets is 0.73--1.00 on ECFP4, 0.44--0.61 on Sort \& Slice and 0.32--0.46 on Avalon. On PDV, ChemBERTa and MHG-GNN it is 0.13--0.22, 0.14--0.23 and 0.21--0.25. Laplace, Student-t and outlier noise give splits like those under Gaussian noise. Grouped-shifted noise raises the epistemic share of both the heteroscedastic GP and QRF, for every dataset and representation. One possible reading is that a shared offset per scaffold family looks to the model like structure in the data, not label noise.

% CODE INPUT
We next asked how closely the aleatoric part follows the amount of added noise. Held-out labels carry no injected noise, so this comparison uses out-of-fold training molecules. We fitted the slope of mean aleatoric uncertainty against the standard deviation of the added noise, both in the label's units, under Gaussian noise. We report the median over folds (\textbf{[FIGURE NEEDED: aleatoric part against the standard deviation of the added noise, per model, ECFP4 and PDV, out-of-fold training molecules, dashed line where reported noise equals added noise]}). On ECFP4 and PDV the slope runs from 0.36, for QRF on PDV on Caco-2, to 1.09, for a full BNN with a variance head on PDV on hERG K$_i$. Including ChemBERTa widens the range to 0.28, for QRF on Caco-2, and 1.62, for a full BNN with a variance head on hERG K$_i$. The aleatoric part can therefore compare noise between datasets or between noise levels. It does not estimate a dataset's own measurement error.

% CODE INPUT
Censoring behaves differently from the other conditions. Clipping labels at a limit narrows their spread. NGBoost and QRF shrink their total uncertainty with that spread on every dataset, and the GPs do the same on hERG K$_i$ and Caco-2. None of these models treats a clipped label as a noisy one. They grow more confident while the labels lose information.

% TODO: MISSING RESULT: the GP uncertainty decomposition by representation.

\subsection{Fingerprint representations carry the usable uncertainty signal}
% [notes moved to thesis_notes_moved.md]

The choice of representation has a much stronger impact on per-sample uncertainty tracking than on noise robustness. Fingerprints and physicochemical descriptors provided a clear signal with which models such as GP and NGBoost could separate noise-induced residuals from structural variation. However, learned embeddings such as MHG-GNN and mol2vec did not produce viable uncertainty-noise correlations (Table~\ref{tab:top_unc_noise}). Although these learned embeddings can be effective representations for many tasks, this research suggests that they provide challenges for per-sample noise tracking. However, there is room for further investigation, as only two embeddings were tested.

\section{Conclusions}
% [notes moved to thesis_notes_moved.md]

These results lead to two main conclusions. The first is that noise robustness, defined as the rate at which predictive performance degrades with increasing label noise, is primarily determined by the model's training mechanism. The margin maximization of SVM, the weight priors of BNNs, and the ensembling in tree methods all appear to keep those models from fitting to noise. The choice of representation has a stronger influence on predictive performance than noise robustness. The interaction term between representation and model has a stronger impact; some models like SVM and full BNNs remain robust across all representations, while RF and NNs are robust only with certain ones. While the pattern of noise has a strong impact on predictive performance, it does not influence noise robustness rankings between models or representations much.

We investigated how a wide variety of molecular representations and model architectures perform under increasing label noise in QSAR modeling. We used ANOVA variance decomposition across three noise conditions on the QM9 HOMO--LUMO gap task to show that with regards to predictive performance, model architecture and molecular representation account for comparable shares of the variance. However, for robustness to label noise, measured as \aucnorm, the normalized area under the $R^2$ retention curve, model architecture is the dominant factor, while representation explains less than 10\% of the variance.

Threshold and value-proportional noise produce the widest spread of NDS across models; outlier noise barely separates them. Model rankings are highly concordant across strategies (Kendall's $W = 0.92$), suggesting that a model's robustness to noise is independent of type of noise. Predictive performance on clean data is not correlated with noise robustness; NGBoost and SVM, the most noise-robust models, did not perform particularly well on clean data. Embeddings (MHG-GNN, mol2vec) struggled with regards to noise robustness, especially with neural network architectures. On the flip side, QRF was significantly less noise-robust than RF across datasets, suggesting that quantile regression may be overfitting to noisy labels rather than absorbing them.

% TODO: review, needs manual review
The second conclusion concerns whether a model's uncertainty estimates track per-sample label noise. Some models struggle, particularly those like BNNs whose uncertainty comes from a posterior over weights or QRF which does produce a full quantile, but is unable to successfully track noise. 
Models that fit a separate, per-sample scale or observation-noise parameter during training are able to track that label noise more effectively. NGBoost identifies noisy labels through a per-sample predicted scale that absorbs large residuals. The GP does not: it has a single global observation noise term, so its per-sample uncertainty comes from the posterior variance, which depends on where a molecule sits relative to the training molecules and not on its label. Among representations, uncertainty tracking collapses with the learned embeddings MHG-GNN and mol2vec, while fingerprints and descriptors produce a viable signal. Neither strong predictive performance nor noise robustness implies noise tracking; again it is up to the model's architecture and how it handles loss. 

Although the choice of representation carries less weight than that of models, we observed an interesting phenomenon. PDVs produce the strongest predictive performance on clean data, but are the least noise-resistant. Embeddings, particularly mol2vec, are the most robust to noise yet perform the worst noise tracking. At the end of the day, the choice of both model and representation comes down to your problem, including the amount of noise in your data, your objectives, and your compute limits. 

QM9 served as a relatively clean baseline; it is computationally-derived with negligible measurement noise. The validation datasets, LogD, Caco-2 \citep{openadmet}, and hERG Ki \citep{Zdrazil2023}, tested are smaller and cover a narrow range of the chemical and biological space, and have an unknown inherent amount of noise from experimentally-determined values, on top of the added artificial noise. These validation datasets supported the QM9 findings, with XGBoost the notable exception, degrading on the external datasets. Future work should extend this type of noise robustness benchmarking to classification tasks and larger experimental datasets with better-known uncertainty estimates.

When working with noisy experimental data in QSAR settings, it is important to keep in mind that noise robustness and noise-aware uncertainty estimation are primarily driven by the model's mechanism of training, and to a certain extent the choice of molecular representation, usually with respect to the model. While Bayesian transformations on neural networks improve noise robustness, they do not produce uncertainty estimates that correlate particularly strongly with noise. Models like NGBoost and GPs, which include a learned noise or scale parameter in the training objective, are the most robust to noise across both various types of noise and datasets and often produce uncertainty estimates which track per-sample label noise. Pairing these models with fingerprint-based representations provides additional benefits in restricting, and potentially detecting and mitigating label noise. 

